\clearpage
\SourceLeaf{237}%
\begin{center}
\Large\textsc{[Lübsens Parallelentheorie.]}
\end{center}


\begin{center}
\emph{[\textup{\textsc{Schumacher}} an \textup{\textsc{Gauss}}. Altona, 31. Dezember 1835.]}
\end{center}

\{\mbox{.\ .\ .\ .\ .\ .\ .\ .\ .} \textsc{Lübsen} quält mich sehr mit seiner Parallelentheorie
(wahrscheinlich greift hier die Nemesis ein, um mir meine eigenen
Versuche zu vergelten). Ich lege Ihnen seinen Beweis bei. Ich habe ihm
vergeblich gesagt, dass sich Linien besser mit dem
Parallellineal als in der Theorie verschieben lassen. Er bleibt aber bei
seiner Überzeugung, den geometrischen Beweis gefunden zu haben.
Da er sonst ein kluger und
bescheidener Mann ist, wäre es vielleicht gut, wenn Sie ihn mit ein paar Worten von seinem
Irrtum abbringen könnten. Es liegt wohl daran, dass ich ihm seinen
Irrtum nicht so deutlich machen kann, wie ich ihn selbst
erkenne. \mbox{.\ .\ .\ .\ .\ .}\}


\begin{center}
\textsc{Gauss} an \textsc{Schumacher}. Göttingen, 2. Januar 1836.
\end{center}

\mbox{.\ .\ .\ .\ .\ .\ .\ .} Der Kern des Fehlers von Herrn \textsc{Lübsen} besteht darin,
dass Euklids Geometrie falsch sein kann, ohne dass es in seiner
Konstruktion bei
